% Ad Familiares 14.2 --- headnote
% ad-familiares-14-02, 5 October 58 BC, written at Thessalonica

\textit{Cicero to Terentia, Tullia, and the boy Marcus, written
from Thessalonica on the third day before the Nones of October
(5 October) 58 BC --- the second surviving letter to the family
from exile, after the Brundisium letter \textit{Fam.} 14.4 of
late April. The opening half is the bearing-up of his own grief
through Tullia's husband Piso, and the calculation about the
new tribunes (Pompey willing, Crassus to be feared still).
\textsection 2 carries one of the most painful incidents of
the surviving correspondence: Cicero has just learned, through
a letter from P. Valerius, that Terentia had been hauled from
the precinct of Vesta to the bankers' Tabula Valeria in the
Forum --- a public humiliation, possibly to make her account
for property forfeit under the Clodian law. The line is
Cicero's: ``you, my light, my heart's longing, the one to whom
all used to look for help.'' \textsection 3 is the tug between
his fear that Terentia will throw the last of her own
\textit{dos} into the cost of recovering the Palatine house,
and his recognition that ``everything is in you.'' He begs her
to spare her health.}
